\begin{solution}
Metric compatibility gives three cyclic equations,
\[
\partial_\mu g_{\nu\sigma}
=\Gamma^\rho{}_{\mu\nu}g_{\rho\sigma}
+\Gamma^\rho{}_{\mu\sigma}g_{\nu\rho}.
\]
Adding the \(\mu,\nu\) equations and subtracting the \(\sigma\) equation, then
using symmetry in the lower connection indices, yields
\[
\Gamma^\rho{}_{\mu\nu}
=\frac12g^{\rho\sigma}
(\partial_\mu g_{\nu\sigma}+\partial_\nu g_{\mu\sigma}
-\partial_\sigma g_{\mu\nu}).
\]
The difference of two compatible torsion-free connections is simultaneously
symmetric in one pair and antisymmetric after lowering an index in another,
forcing it to vanish.
\end{solution}

\begin{solution}
Using
\(\Gamma^\sigma{}_{\rho\sigma}=\partial_\rho\ln\sqrt{|g|}\), the derivative
of \(\sqrt{|g|}\) cancels the sum of connection terms acting on the \(n\)
indices of the alternating symbol.  Hence \(\nabla_\rho\varepsilon=0\).
Because the Hodge star is built only from \(g^{\mu\nu}\) and
\(\varepsilon_{\mu_1\ldots\mu_n}\), both covariantly constant,
\(\nabla(\star\omega)=\star(\nabla\omega)\).
\end{solution}

\begin{solution}
The identity
\[
\sqrt{|g|}\,\nabla_\mu V^\mu
=\partial_\mu(\sqrt{|g|}V^\mu)
\]
reduces the result to the ordinary divergence theorem in coordinate patches;
partition of unity makes it global.  The directed surface element is
\(n_\mu\dd\Sigma\), with outward spacelike normal on timelike boundary pieces
and future/past-directed timelike normals on final/initial spacelike pieces.
\end{solution}

\begin{solution}
If \(u=f\,\dd t\), then \(u\wedge\dd u=f\,\dd t\wedge\dd f\wedge\dd t=0\).
Conversely, Frobenius' theorem says the kernel distribution of \(u\) is
integrable precisely when \(u\wedge\dd u=0\), producing local level surfaces
\(t=\) constant and hence \(u=f\,\dd t\).  In components,
\[
u_{[\mu}\nabla_\nu u_{\rho]}=0.
\]
For normalized timelike \(u^\mu\), this is equivalent to vanishing projected
vorticity.
\end{solution}

\begin{solution}
The first Cartan equation gives
\[
\dd e^1+\omega^1{}_2\wedge e^2=0,\qquad
\dd e^2+\omega^2{}_1\wedge e^1=0,
\]
so \(\omega^1{}_2=-\dd\phi\) and
\(\omega^2{}_1=\dd\phi\).  The curvature is
\[
\Omega^1{}_2=\dd\omega^1{}_2
+\omega^1{}_a\wedge\omega^a{}_2=0
\]
away from the coordinate origin.  The connection records the rotation of the
polar frame, not intrinsic curvature.
\end{solution}

\begin{solution}
Direct substitution gives
\[
\widetilde\Gamma^\rho{}_{\mu\nu}
-\Gamma^\rho{}_{\mu\nu}
=\delta^\rho_\mu\partial_\nu\omega
+\delta^\rho_\nu\partial_\mu\omega
-g_{\mu\nu}\partial^\rho\omega.
\]
For a null tangent \(k^\mu\), the added term in the geodesic equation is
\(2(k\cdot\partial\omega)k^\rho\); the term proportional to \(k^2\) vanishes.
It is tangent to the curve and can be absorbed by reparameterization.
\end{solution}

\begin{solution}
Normal coordinates are defined by the exponential map, so radial geodesics
are straight and
\(\Gamma^\rho{}_{(\mu\nu,\alpha)}(p)=0\), while
\(\Gamma^\rho{}_{\mu\nu}(p)=0\).  Combining these identities with the
definition of curvature and metric compatibility gives
\[
\partial_\alpha\partial_\beta g_{\mu\nu}(p)
=-\frac13\bigl(R_{\mu\alpha\nu\beta}
+R_{\mu\beta\nu\alpha}\bigr).
\]
Taylor expansion yields the stated result.  The quadratic coefficient is a
tensorial curvature invariant and therefore cannot be removed by coordinates.
\end{solution}

\begin{solution}
Expanding each flow to first order and retaining the mixed term gives, on any
test function \(f\),
\[
f\longmapsto f+st(XY-YX)f+\order(s^2t,st^2)
=f+st[X,Y]f+\cdots.
\]
Thus the infinitesimal parallelogram made from the two flows fails to close by
the bracket vector.  Vanishing bracket is precisely the local commutativity of
the flows.
\end{solution}
